\documentclass[12pt,a4paper]{article}
\usepackage{graphicx} % Required for inserting images

\title{LaTeX Report}
\author{Syed Muhammad Shamlan Ali}
\date{September 2026}

\begin{document}

\maketitle

\section{Introduction}
The portfolio I created was to prove technical coding skills as well as demonstrate personal projects to potential employers. It's almost like a personal exhibit of your best work that helps with getting future
positions at different companies. I created this portfolio based on a Computer Science final project, and it's 
supposed to show very brief beginner-level coding skills. It demonstrates all the knowledge and skills that I have achieved and learned so far.


\section{Portfolio Overview}
The Portfolio contains a bit of personal information about me, as well as my interests and contact information. It also displays all the skills I've gained throughout my academic journey and shows a brief project I've created over time, which is a Python calculator. It basically collects two numbers from the user and an operator and calculates the result. My CV is also attached to the portfolio to give future employers a brief explanation of my work and skills.

\section{Portfolio Web Structure}
I think it's valuable to give a short, brief explanation about yourself to give people an idea of what kind of person you might be, which is beneficial to how your work ethic. I briefly talk about my academic journey and my interests. I mention some of my technical skills, such as working with Python and LaTeX, and having soft skills
such as teamwork, willingness to learn, and critical thinking. On the same portfolio, visitors or future employers can access my CV by clicking the link. My contact info is there, and they can communicate me through my Gmail.

\section{GitHub Repository Structure}
I organized my GitHub repository into files. I created a folder called Projects, then inserted my Python calculator's source code and a screenshot of it working in there, with a brief explanation of how it works. Then I created another folder called CV, which has my CV file and the source code of the CV, with a brief explanation of what they are too.

\section{Design Decisions}
When creating the portfolio, this was my first time doing something like this for a Computer science course, so
I went for the very basic route. I chose a black background for a more professional and average look, and the Times New Roman font since it's the font used for formatting documents. I chose size 12 so it's readable for anyone viewing the portfolio. By clicking on the words on top, they clearly navigate you to those sections, which is pretty simple. That was kinda my main thought process behind the design of the website. I was just looking for a simple, normal design, but later on, when I got to improving it, I went for a slightly flashier look.

\section{Tools and Technologies}
I used LaTeX to create the CV Format and the final report. GitHub Repository was used to collect all the files in one spot where they could all be easily accessible. For creating the Portfolio, I used HTML for the code of the portfolio and CSS for the design, and by using GitHub pages it created a portfolio.

\section{Implementation}
I started with the CV first since it was the fastest task to be completed; I used LaTeX for the formatting. I then
created the GitHub Repository and then added the folder for the CV with the source code and CV. I then created
a quick little Python Project, which was the calculator, as I needed something to show off my coding skills. I then created a folder for that on the GitHub Repository and uploaded all the necessary details. My final task was creating the Portfolio by creating the HTML and CSS on GitHub, where I used GitHub Pages to connect them. And finally created my LaTeX report.

\section{Strengths and Weaknesses}
When creating the portfolio, I believe the easiest part was creating the sections and creating a navigation system for the portfolio. They didn't require much knowledge, and I was able to attach the CV and my Python calculator. Though I believe designing the portfolio is a bit more difficult than I thought. I was only able to give it a bland, average look because I wasn't very knowledgeable, and that limited my creativity.

\section{Future Improvements}
I think by far the biggest improvement is the design of the portfolio, as it is very lacking in creativity, and my projects. I never really worked on many projects for computer science, so the only thing I had to show was a very brief and beginner-level Python calculator. As I create and improve the design of the portfolio, I would need to also work on more personal projects to stand out more to future employers and visitors.


\section{Conclusion}
Throughout this project journey, I was successfully able to create a GitHub repository and attach all my work and findings there in an organized structure. My projects and CV, with their source code and brief explanations, are safely organized in the folder. I was able to use HTML and CSS to create a website portfolio that is able to display my work and skills by using GitHub Pages. I also gained a brief and clear understanding of how to create a document using LaTeX formatting.


\section{References}
Git Hub (2026) \textit{GitHub Pages documentation}. Available at: \url{https://docs.github.com/en/pages}(Accessed: 12 September 2026)

MDN Web Docs (2026) \textit{HTML: HyperText Markup Language}. Available at: \url{https://developer.mozilla.org/en-US/docs/Web/HTML}(Accessed: 15 September 2026)

MDN Web Docs (2026) \textit{Cascading Style Sheets}. Available at: \url{https://developer.mozilla.org/en-US/docs/Web/CSS}(Accessed: 15 September 2026)

MDN Web Docs (2026) \textit{ Getting Started with CSS}. Available at: \url{https://developer.mozilla.org/en-US/docs/Learn_web_development/Core/Styling_basics/Getting_started}(Accessed: 15 September 2026)



\end{document}
